\listfiles
\documentclass{slides}
\usepackage[screen,includeheadfoot,left=1em,right=1em,top=0.5em,bottom=0.2em]{geometry}

\usepackage{iftex}
\RequireLuaTeX

% \usepackage{fontspec}
\usepackage{fancyhdr}
\usepackage{enumitem}
\usepackage{xcolor}
\usepackage{listings}
\usepackage{tcolorbox}
\usepackage{fontawesome}
\usepackage[
  hyperfootnotes=false,
  pdfpagelabels=false,
  linktoc=all,
  colorlinks=true,
  linkcolor=black,
  citecolor=black,
  urlcolor=cern80,
  pdfusetitle,
  pdfsubject={},
  pdfkeywords={}
]{hyperref}

% Fonts ------------------------------------------------------------------------
\setsansfont{Roboto}[
  Path = fonts/Roboto/,
  Extension = .ttf,
  UprightFont = *-Regular,
  ItalicFont = *-Italic,
  BoldFont = *-Bold,
  BoldItalicFont = *-BoldItalic,
  FontFace = {l}{n}{*-Light},
  FontFace = {l}{it}{*-LightItalic},
  FontFace = {sb}{n}{*-Medium},
  FontFace = {sb}{it}{*-MediumItalic},
  FontFace = {eb}{n}{*-Black},
  FontFace = {eb}{it}{*-BlackItalic}
]

\setmonofont{DejaVu Sans Mono}

\NewDocumentCommand\F{mO{#1}}{\fontsize{#1}{#2}\selectfont}

% Colors -----------------------------------------------------------------------
\definecolor{cern}{HTML}{0033A0}
\definecolor{cern80}{HTML}{385CB4}

\definecolor{araBlue}{HTML}{003376}
\definecolor{araBlue2}{HTML}{0092CA}
\definecolor{araOrange}{HTML}{DE5700}

\definecolor{f6}{HTML}{F6F6F6}
\definecolor{green80}{HTML}{008000}

% Headers and Footers ----------------------------------------------------------
\fancyhf{} % Clear all default header and footer fields

\makeatletter
\newcommand{\@slidetitle}{}
\newcommand{\slidetitle}[1]{\renewcommand{\@slidetitle}{#1}}

\providecommand{\lastpage}{0}
\AtEndDocument{
  \write\@auxout{\string\gdef\string\lastpage{\thepage}}
}

\lhead{\fontseries{sb}\selectfont\Large\color{araBlue}\@slidetitle}
\lfoot{\F{11pt}\color{araBlue}%
  \hypersetup{hidelinks}\href{https://github.com/ivankp/slides-nlohmann-json}{\@title}}
\cfoot{\F{11pt}\color{araBlue}\@author}
\rfoot{\F{11pt}\color{araBlue}\@date
  \begin{minipage}[b]{5em}%
    \hfill%
    {\F{14pt}\thepage\nobreak\hspace{1pt}}/\nobreak\hspace{1pt}\lastpage
  \end{minipage}%
}
\makeatother

\fancyfootoffset{1em}

% Itemize, Enumerate -----------------------------------------------------------
\setlist{
    topsep=0pt,
    parsep=0em,
    partopsep=0pt,
    itemsep=0pt,
    leftmargin=1em
}

\setlist[1]{itemsep=1em}
\setlist[2]{topsep=0.25em, itemsep=0.25em, before=\F{18pt}[24pt]}

\setlist[itemize,1]{label={\color{araOrange}$\bullet$}}
\setlist[itemize,2]{label={\color{araOrange}$\circ$}}
\setlist[itemize,3]{label={\color{araOrange}\bfseries --}}

% Colorbox ----------------------------------------------------------------------------------------
\tcbset{
  colback=f6,
  colframe=araBlue
}

% Listings ---------------------------------------------------------------------
\definecolor{lstType}{HTML}{008000} % cppreference k
\definecolor{lstKeyword}{HTML}{440099}
\definecolor{lstLiteral}{HTML}{B00040} % cppreference kt
\definecolor{lstNamespace}{HTML}{505050}
\definecolor{lstFunction}{HTML}{0077AA}
\definecolor{lstComment}{HTML}{708090}
\definecolor{lstPreproc}{HTML}{9C6500} % cppreference cp

\lstset{
  aboveskip=0pt,
  belowskip=0pt,
%
  xleftmargin=\parindent,
%
  upquote=true,
%
  basicstyle={\ttfamily},
%
  escapeinside={@!}{@},
%
  moredelim=[is][\color{lstType}]{@T}{@},
  moredelim=[is][\bfseries\color{lstType}]{@bT}{@},
  moredelim=[is][\color{lstKeyword}]{@K}{@},
  moredelim=[is][\color{lstLiteral}]{@L}{@},
  moredelim=[is][\color{lstNamespace}]{@N}{@},
  moredelim=[is][\color{lstFunction}]{@F}{@},
  moredelim=[is][\color{lstPreproc}]{@P}{@},
% moredelim=[l][\color{lstPreproc}]{\#},
%
  morecomment=[l]{//}, % l is for line comment
  morecomment=[s]{/*}{*/}, % s is for start and end delimiter
  commentstyle={\itshape\color{lstComment}},
%
  morestring=[b]", % "string"
  morestring=[b]', % 'c'
  morestring=[s]{R"(}{)"}, % R"(raw string)"
  stringstyle=\color{lstLiteral},
  showstringspaces=false,
%
  morekeywords=[0]{
alignas, alignof, and, and_eq, asm, bitand, bitor, break, case, catch, class, co_await, co_return,
co_yield, compl, concept, const, const_cast, consteval, constexpr, constinit, continue,
contract_assert, decltype, default, delete, do, dynamic_cast, else, enum, explicit, export, extern,
for, friend, goto, if, inline, mutable, namespace, new, noexcept, not, not_eq, operator, or, or_eq,
private, protected, public, reflexpr, register, reinterpret_cast, requires, return, sizeof, static,
static_assert, static_cast, struct, switch, synchronized, template, this, thread_local, throw, try,
typedef, typeid, typename, union, using, virtual, volatile, while, xor, xor_eq,
nullptr
  },
  keywordstyle=[0]\color{lstKeyword},
%
  morekeywords=[1]{
auto, void, bool, char, int, short, long, signed, unsigned, float, double,
int8_t, uint8_t, int16_t, uint16_t, int32_t, uint32_t, int64_t, uint64_t,
size_t
},
  keywordstyle=[1]\color{lstType},
%
  morekeywords=[2]{
true, false
},
  keywordstyle=[2]\color{lstLiteral},
%
  literate=
    {::}{{{\color{lstNamespace}{::}}}}{2}
    {std::}{{{\color{lstNamespace}{std::}}}}{5}
}

\definecolor{jsep}{HTML}{670026}

\lstdefinelanguage{json}{
  basicstyle={\ttfamily\color{lstLiteral}},
  upquote=true,
  showstringspaces=false,
  literate=
    {:}{{{\color{jsep}{:}}}}{1}
    {,}{{{\color{jsep}{,}}}}{1}
    {\{}{{{\color{jsep}{\{}}}}{1}
    {\}}{{{\color{jsep}{\}}}}}{1}
    {[}{{{\color{jsep}{[}}}}{1}
    {]}{{{\color{jsep}{]}}}}{1},
}

% Title, Authors, Date *********************************************************
\title{\texorpdfstring{\ttfamily nlohmann/json}{nlohmann/json}}
\author{Ivan Pogrebnyak}
\date{\today}
% \date{April 21, 2026}

%%%%%%%%%%%%%%%%%%%%%%%%%%%%%%%%%%%%%%%%%%%%%%%%%%%%%%%%%%%%%%%%%%%%%%%%%%%%%%%%
\begin{document}

\thispagestyle{empty}
\makeatletter
{\centering \huge \color{araBlue} \@title \par}
{\centering \Large \color{araBlue2} {JSON for Modern C++} \par}
{\centering \Large \@author \par}
{\centering \@date \par}
\makeatother
\pagebreak
\pagestyle{fancy}

% ------------------------------------------------------------------------------
\slidetitle{JSON}

\begin{itemize}
\item JavaScript Object Notation
\item Data types:
  null, boolean, number, string, array, object
\item {\bfseries Tree structure}: branches are either objects or arrays
  \begin{itemize}
    \item Object keys are strings
    \item Array keys are indices
  \end{itemize}
\item No comments
\end{itemize}

\begin{tcolorbox}
\begin{lstlisting}[language=json]
{
    "key": [
        null, true, false, 42, 1.41, "string"
    ]
}
\end{lstlisting}
\end{tcolorbox}

\pagebreak
% ------------------------------------------------------------------------------
\slidetitle{\lstinline{nlohmann/json}}

\begin{itemize}[itemsep=0.5em]
  \item \faGithub\ \ \url{https://github.com/nlohmann/json}
  \item \faFileTextO\ \ \url{https://json.nlohmann.me/api/basic_json/}
  \begin{itemize}
    \item \faFileTextO\ \ in the slides link to documentation
  \end{itemize}
  \item Developed since 2013
  \item Header-only library
  \item {\color{orange}\faWarning}\ \ Not a fast library but easy to use
  \item {\color{orange}\faWarning}\ \ Parses whole document
  \item {\color{orange}\faWarning}\ \ Handles errors via exception
    \begin{itemize}
      \item Unless you explicitly use non-throwing interfaces
    \end{itemize}
  \item Similar in spirit to the standard library container templates
  \item Uses standard library types to represent JSON values
    \begin{itemize}
      \item This is configurable via \lstinline{@Tbasic_json@} template parameters
    \end{itemize}
\end{itemize}

\pagebreak
% ------------------------------------------------------------------------------
\slidetitle{\ttfamily nlohmann/json}

\begin{tcolorbox}
{\F{12pt}[16pt]
\begin{lstlisting}
namespace @Nnlohmann@ {

template <
    template <typename, typename, typename...> class @bTObjectType@ = std::@Tmap@,
    template <typename, typename...> class @bTArrayType@ = std::@Tvector@,
    class @bTStringType@ = std::@Tstring@,
    class @bTBooleanType@ = bool,
    class @bTNumberIntegerType@ = std::int64_t,
    class @bTNumberUnsignedType@ = std::uint64_t,
    class @bTNumberFloatType@ = double,
    template <typename> class @bTAllocatorType@ = std::@Tallocator@,
    template <typename, typename = void> class @bTJSONSerializer@ = @Tadl_serializer@,
    class @bTBinaryType@ = std::@Tvector@<std::uint8_t>,
    class @bTCustomBaseClass@ = void
>
class @Tbasic_json@;

using @Tjson@ = @Tbasic_json@<>;
using @Tordered_json@ = @Tbasic_json@<@Tordered_map@>;

template <class @bTKey@, class @bTT@>
struct @Tordered_map@ : std::@Tvector@<std::@Tpair@<const @bTKey@, @bTT@>>;

} // namespace nlohmann
\end{lstlisting}
}
\end{tcolorbox}

\pagebreak
% ------------------------------------------------------------------------------
\slidetitle{Parsing}

\begin{tcolorbox}
\begin{lstlisting}
static @Tbasic_json@ @Tbasic_json@::@Fparse@(); @!{\hfill
  \href{https://json.nlohmann.me/api/basic_json/parse/}{\faFileTextO}}@
\end{lstlisting}
\end{tcolorbox}
\vspace{0em}

\begin{itemize}[itemsep=0.5em]
  \item Valid arguments:
    \begin{itemize}
      \item \lstinline{std::@Tistream@}
      \item \lstinline{@TFILE@*}
      \item null-terminated \lstinline{char*} or \lstinline{uint8_t*}
      \item an object \lstinline{obj} with \lstinline{@Fbegin@(obj)} and
        \lstinline{@Fend@(obj)}: \lstinline{std::@Tstring@}
      \item a pair of iterators
    \end{itemize}

  \item Examples:
    \begin{itemize}
      \item Parse a file
        \vspace{-0.5em}
        \begin{tcolorbox}{\F{16pt}
        \begin{lstlisting}
const @Tjson@ j @F=@ @Tjson@::@Fparse@(std::@Tifstream@("data.json"));
        \end{lstlisting}
        }\end{tcolorbox}

      \item Parse a string
        \vspace{-0.5em}
        \begin{tcolorbox}{\F{16pt}
        \begin{lstlisting}
const @Tjson@ j @F=@ @Tjson@::@Fparse@(R"(@!\color{jsep}\{@ "key"@!\color{jsep}:@ "value" @!\color{jsep}\}@)");
        \end{lstlisting}
        }\end{tcolorbox}
    \end{itemize}
\end{itemize}
\vspace{-1em}

\pagebreak
% ------------------------------------------------------------------------------
\slidetitle{Parsing}

\begin{tcolorbox}
\begin{lstlisting}
std::@Tistream@&
operator@F>>@(std::@Tistream@&, @Tbasic_json@&); @!\hfill
  \href{https://json.nlohmann.me/api/operator_gtgt/}{\faFileTextO}@
\end{lstlisting}
\end{tcolorbox}

\begin{itemize}
  \item Can read json data from an input stream
    \begin{tcolorbox}{\F{16pt}[26pt]
    \begin{lstlisting}
@Tjson@ j; // can't be const
std::@Tifstream@("/opt/data.json") @F>>@ j;
    \end{lstlisting}
    }\end{tcolorbox}
\end{itemize}

\pagebreak
% ------------------------------------------------------------------------------
\slidetitle{Parsing}

\begin{tcolorbox}
\begin{lstlisting}
static bool @Tbasic_json@::@Faccept@(); @!\hfill
  \href{https://json.nlohmann.me/api/basic_json/accept/}{\faFileTextO}@
\end{lstlisting}
\end{tcolorbox}

\begin{itemize}
  \item Accepts same arguments as \lstinline{@Fparse@()}
  \item Only checks if the input is valid JSON
\end{itemize}

\pagebreak
% ------------------------------------------------------------------------------
\slidetitle{Parsing}

\begin{tcolorbox}{\F{18pt}[20pt]
\begin{lstlisting}
@Tjson@ operator ""@F_json@(const char*, std::size_t); @!\hfill
  \href{https://json.nlohmann.me/api/operator_literal_json/}{\faFileTextO}@
\end{lstlisting}
}\end{tcolorbox}

\begin{itemize}
  \item User-defined literal operator
  \item {\color{orange}\faWarning}\ \ Parses at run time
  \item Not \lstinline{constexpr}
  \item Example:
    \begin{tcolorbox}{\F{16pt}[26pt]
    \begin{lstlisting}
using namespace @Nnlohmann@::@Nliterals@;
const @Tjson@ j @F=@ R"(@!\color{jsep}\{@ "key"@!\color{jsep}:@ "value" @!\color{jsep}\}@)"@F_json@;
    \end{lstlisting}
    }\end{tcolorbox}
\end{itemize}

\pagebreak
% ------------------------------------------------------------------------------
\slidetitle{Emitting}

\begin{tcolorbox}{\F{19pt}
\begin{lstlisting}
@Tbasic_json@::@Tstring_t@ @Tbasic_json@::@Fdump@() const; @!\hfill
  \href{https://json.nlohmann.me/api/basic_json/dump/}{\faFileTextO}@
\end{lstlisting}
}\end{tcolorbox}

\begin{itemize}
  \item Non-static member function
  \item Writes to a string
  \begin{itemize}
    \item {\color{orange}\faWarning}\ \ Use \lstinline{operator@F<<@}{}
      to write directly to a file
  \end{itemize}
  \item Example:
    \begin{tcolorbox}{\F{16pt}[26pt]
    \begin{lstlisting}
const @Tjson@ j;
const std::@Tstring@ str @F=@ j.@Fdump@();
    \end{lstlisting}
    }\end{tcolorbox}
\end{itemize}

\pagebreak
% ------------------------------------------------------------------------------
\slidetitle{Emitting}

\begin{tcolorbox}{\F{19pt}[28pt]
\begin{lstlisting}
std::@Tostream@&
operator@F<<@(std::@Tostream@&, const @Tbasic_json@&); @!\hfill
  \href{https://json.nlohmann.me/api/operator_ltlt/}{\faFileTextO}@
\end{lstlisting}
}\end{tcolorbox}

\begin{itemize}
  \item Can directly write to a file
    \begin{tcolorbox}{\F{16pt}[26pt]
    \begin{lstlisting}
const @Tjson@ j;
std::@Tofstream@("/opt/data.json") @F<<@ j;
    \end{lstlisting}
    }\end{tcolorbox}
\end{itemize}

\pagebreak
% ------------------------------------------------------------------------------
\slidetitle{Exceptions}

\begin{itemize}
  \item \begin{lstlisting}
@Nnlohmann@::@Tbasic_json@::@Texception@; @!\hfill
      \href{https://json.nlohmann.me/api/basic_json/exception/}{\faFileTextO}@
    \end{lstlisting}
  \item \lstinline{class @Texception@ : public std::@Texception@;}
  \item Sub-classes:
    \begin{itemize}
      \item \lstinline{@Tparse_error@}:
        a parsing error
      \item \lstinline{@Tout_of_range@}:
        access out of the defined range
      \item \lstinline{@Ttype_error@}:
        executing a member function with a wrong type
      \item \lstinline{@Tinvalid_iterator@}:
        errors with iterators
      \item \lstinline{@Tother_error@}:
        other library errors
    \end{itemize}
  \item Full list of exceptions \hfill
    \href{https://json.nlohmann.me/home/exceptions/}{\faFileTextO}
\end{itemize}

\pagebreak
% ------------------------------------------------------------------------------
\slidetitle{\faHandORight\,\ Best practice: Use exceptions}

\begin{itemize}
  \item Exceptions provide diagnostic information.
  \item Most {\ttfamily nlohmann} interfaces don't have non-throwing options.
  \item {\ttfamily nlohmann} interfaces that can be used without exceptions
    can only report the fact of failure with no clarifying information.
    \begin{itemize}
      \item \lstinline{@Fparse@()} returns \lstinline{@Tbasic_json@} in
        \lstinline{discarded} state.
    \end{itemize}
  \item Exceptions can be switched off for the whole library \hfill
    \href{https://json.nlohmann.me/home/exceptions/#switch-off-exceptions}{\faFileTextO}
    \begin{itemize}
      \item \lstinline{@Fabort@()} is called anywhere the library would
        \lstinline{throw}.
    \end{itemize}
\end{itemize}

\pagebreak
% ------------------------------------------------------------------------------
\slidetitle{Exceptions examples}

% \begin{tcolorbox}
\begin{lstlisting}
const @Tjson@ j @F=@ R"(@!\color{jsep}\{@ "a"@!\color{jsep}:@ "value" @!\color{jsep}\}@)"@F_json@;
\end{lstlisting}
% \end{tcolorbox}

\begin{itemize}
  \item \begin{lstlisting}
auto& b @F=@ j.@Fat@("b"); // auto @!{\color{lstComment}$\rightarrow$}@ const json
    \end{lstlisting}
    \begin{itemize}
      \item \color{red}\verb|[json.exception.out_of_range.403]| \\
        \verb|key 'b' not found|
    \end{itemize}

  \item \begin{lstlisting}
int a @F=@ j.@Fat@("a");
    \end{lstlisting}
    \begin{itemize}
      \item \color{red}\verb|[json.exception.type_error.302]| \\
        \verb|type must be number, but is string|
    \end{itemize}
\end{itemize}

\pagebreak
% ------------------------------------------------------------------------------
\slidetitle{\lstinline{JSON_DIAGNOSTICS}}

\begin{lstlisting}
@P#define JSON_DIAGNOSTICS@ @L1@ @!\hfill
      \href{https://json.nlohmann.me/api/macros/json_diagnostics/}{\faFileTextO}@
@P#include@ @L<nlohmann/json.hpp>@@!\vspace{-0.8em}@

const @Tjson@ j @F=@ R"(@!\color{jsep}\{@ "a"@!\color{jsep}:@ @!\color{jsep}\{@ "b"@!\color{jsep}:@ 1 @!\color{jsep}\}@ @!\color{jsep}\}@)"@F_json@;
\end{lstlisting}
\vspace{-0.5em}

\begin{itemize}[itemsep=0.2em]
  \item \begin{lstlisting}
auto& c @F=@ j.@Fat@("a").@Fat@("c");
    \end{lstlisting}
    \begin{itemize}
      \item \color{red}\verb|[json.exception.out_of_range.403]| \\
        \verb|(/a) key 'c' not found|
    \end{itemize}

  \item \begin{lstlisting}
int a @F=@ j.@Fat@("a");
    \end{lstlisting}
    \begin{itemize}
      \item \color{red}\verb|[json.exception.type_error.302]| \\
        \verb|(/a) type must be number, but is object|
    \end{itemize}

  \item {\F{18pt} When enabled, exception messages contain a JSON Pointer to
    the JSON value that triggered the exception.}

  \item {\F{18pt} {\color{orange}\faWarning}\ \ Note that enabling this macro increases
    the size of every JSON value by one pointer and adds some runtime overhead.}
\end{itemize}

\pagebreak
% ------------------------------------------------------------------------------
\slidetitle{JSON node access}

\begin{itemize}[itemsep=0.6em]
  \item \lstinline{operator@F[]@()} models mutable access
    \begin{itemize}
      \item {\color{orange}\faWarning}\ \ Inserts missing keys if the object is
        not \lstinline{const}
      \item Throws on missing keys if the object is \lstinline{const}
    \end{itemize}
  \item \lstinline{@Fat@()} models immutable access
    \begin{itemize}
      \item Always throws on missing keys
    \end{itemize}
  \item Both return references to the \lstinline{@Tvalue_type@}
    \begin{itemize}
      \item i.e. \lstinline{@Nnlohmann@::@Tbasic_json@&}
    \end{itemize}
  \item What happens if key does not exist:
\end{itemize}
\vspace{-0.5em}

{\F{15.75pt}[24pt]
\begin{tabular}{llll}
\hline
&
\lstinline{std::@Tvector@} &
\lstinline{std::@Tmap@} &
\lstinline{@Nnlohmann@::@Tjson@} \\
\hline
\lstinline{operator@F[]@()} & undefined behavior & insert & insert \hfill
  \href{https://json.nlohmann.me/api/basic_json/operator%5B%5D/}{\faFileTextO} \\
\lstinline{@Fat@()} & throw & throw & throw \hfill
  \href{https://json.nlohmann.me/api/basic_json/at/}{\faFileTextO} \\
\lstinline{operator@F[]@() const} & undefined behavior & doesn't exist & throw \\
\lstinline{@Fat@() const} & throw & throw & throw \\
\hline
\end{tabular}
}

\pagebreak
% ------------------------------------------------------------------------------
\slidetitle{\faHandORight\,\ Best practice: \lstinline{json.at()}}

\begin{itemize}
  \item Prefer \lstinline{@Fat@()} to \lstinline{operator@F[]@()}
    for working with json objects that you don't intend to modify
    to prevent accidental insertion of default values.
\end{itemize}

\vspace{1em}
\begin{lstlisting}
@Tjson@ j @F=@ R"(@!\color{jsep}\{@ "a"@!\color{jsep}:@ "value" @!\color{jsep}\}@)"@F_json@;
\end{lstlisting}
\vspace{-0.5em}

\begin{itemize}[itemsep=0.5em]
  \item \begin{lstlisting}
auto& b @F=@ j.@Fat@("b");
    \end{lstlisting}
    \begin{itemize}
      \item Throws since \lstinline{"b"} doesn't exist
    \end{itemize}

  \item \begin{lstlisting}
auto& b @F=@ j["b"];
    \end{lstlisting}
    \begin{itemize}
      \item Inserts \lstinline{@Lnull@} at key \lstinline{"b"}
      \item Returns a reference to the inserted \lstinline{@Lnull@}
    \end{itemize}

  \item \begin{lstlisting}
auto& b2 @F=@ j["b1"]["b2"];
    \end{lstlisting}
    \begin{itemize}
      \item {\F{16pt}\begin{lstlisting}
@!{\F{20pt}j\,\sffamily{} is now\ \ }@@L@!\color{jsep}\{@ "a"@!\color{jsep}:@ "value", "b1"@!\color{jsep}:@ @!\color{jsep}\{@ "b2"@!\color{jsep}:@ null @!\color{jsep}\}@ @!\color{jsep}\}@@
        \end{lstlisting}}
      \item {\color{orange}\faWarning}\ \ \lstinline{operator@F[]@} can turn
        \lstinline{@Lnull@} into an \lstinline{@Lobject@}
    \end{itemize}
\end{itemize}


\pagebreak
% ------------------------------------------------------------------------------
\slidetitle{Value access}

\begin{tcolorbox}
\begin{lstlisting}
template <typename @TValue@>
@PJSON_EXPLICIT@
@Tbasic_json@::operator @TValue@() const; @!\hfill
  \href{https://json.nlohmann.me/api/basic_json/operator_ValueType/}{\faFileTextO}@
\end{lstlisting}
\end{tcolorbox}
\vspace{0em}

\begin{itemize}[itemsep=0.7em]
  \item {\bfseries Implicit} cast operator
  \item {\color{red}\faWarning}\ \ Can silently converts between numeric types
    \begin{itemize}
      \item Can explicitly check with \lstinline{@Fis_number_integer@()} etc. \hfill
        \href{https://json.nlohmann.me/features/types/number_handling/#determine-number-types}{\faFileTextO}
    \end{itemize}
  \item Implemented by calling \lstinline{@Fget@<@TValue@>()}
  \item Returns by value
\end{itemize}

\begin{tcolorbox}
\begin{lstlisting}
@Tjson@ j;
std::@Tstring@ str @F=@ j;
\end{lstlisting}
\end{tcolorbox}

\pagebreak
% ------------------------------------------------------------------------------
\slidetitle{Value access}

\begin{tcolorbox}
\begin{lstlisting}
template <typename @TValue@>
@TValue@ @Tbasic_json@::@Fget@() const; @!\hfill
  \href{https://json.nlohmann.me/api/basic_json/get/}{\faFileTextO}@
\end{lstlisting}
\end{tcolorbox}
\vspace{0em}

\begin{itemize}[itemsep=0.7em]
  \item {\color{red}\faWarning}\ \ Can silently converts between numeric types.
  \item Returns by value. \lstinline{@TValue@} can't be a reference type.
  \item Implementation equivalent to \vspace{-0.5em}
\begin{tcolorbox}
\begin{lstlisting}
@TValue@ tmp; // creates a temporary
@TJSONSerializer@<@TValue@>::@Ffrom_json@(*this, tmp);
return tmp;
\end{lstlisting}
\end{tcolorbox}
\end{itemize}

\begin{tcolorbox}
\begin{lstlisting}
@Tjson@ j;
std::@Tstring@ str @F=@ j.@Fget@<std::@Tstring@>();
\end{lstlisting}
\end{tcolorbox}

\pagebreak
% ------------------------------------------------------------------------------
\slidetitle{Value access}

\begin{tcolorbox}
\begin{lstlisting}
template <typename @TValue@>
@TValue@& @Tbasic_json@::@Fget_to@(@TValue@& v) const; @!\hfill
  \href{https://json.nlohmann.me/api/basic_json/get_to/}{\faFileTextO}@
\end{lstlisting}
\end{tcolorbox}
\vspace{0em}

\begin{itemize}[itemsep=0.7em]
  \item {\color{red}\faWarning}\ \ Can silently convert float to integer and vice versa
  \item Avoids a temporary variable
  \item Implementation equivalent to \vspace{-0.5em}
\begin{tcolorbox}
\begin{lstlisting}
@TJSONSerializer@<@TValue@>::@Ffrom_json@(*this, v);
\end{lstlisting}
\end{tcolorbox}
\end{itemize}

\begin{tcolorbox}
\begin{lstlisting}
@Tjson@ j;
std::@Tstring@ str;
j.@Fget_to@(str);
\end{lstlisting}
\end{tcolorbox}

\pagebreak
% ------------------------------------------------------------------------------
\slidetitle{Value access}

\begin{tcolorbox}
\begin{lstlisting}
template <typename @TReference@>
@TReference@ @Tbasic_json@::@Fget_ref@() @!\itshape\color{lstKeyword}const@; @!\hfill
  \href{https://json.nlohmann.me/api/basic_json/get_ref/}{\faFileTextO}@
\end{lstlisting}
\end{tcolorbox}
\vspace{0em}

\begin{itemize}[itemsep=0.7em]
  \item Returns a reference to the internally stored value.
    \begin{itemize}
      \item No copies are made.
      \item This reference can be mutable.
    \end{itemize}
  \item \lstinline{@TReference@} must be a reference type.
  \item \lstinline{@TReference@} must be \lstinline{const},
    if the object is \lstinline{const}.
  \item {\color{green80}\faThumbsUp}\ \ Throws if types don't match.
    No accidental conversion.
\end{itemize}

\begin{tcolorbox}{\F{19pt}[26pt]
\begin{lstlisting}
@Tjson@ j;
std::@Tstring@& str @F=@ j.@Fget_ref@<std::@Tstring@&>(str);
\end{lstlisting}
}\end{tcolorbox}

\pagebreak
% ------------------------------------------------------------------------------
\slidetitle{Value access}

\begin{tcolorbox}{\F{19pt}[26pt]
\begin{lstlisting}
template <typename @TPointer@>
@TPointer@ @Tbasic_json@::@Fget_ptr@() @!\itshape\color{lstKeyword}const@ noexcept; @!\hfill
  \href{https://json.nlohmann.me/api/basic_json/get_ptr/}{\faFileTextO}@
\end{lstlisting}
}\end{tcolorbox}
\vspace{0em}

\begin{itemize}[itemsep=0.7em]
  \item Like \lstinline{@Fget_ref@()}, but with pointers.
  \item Returns \lstinline{nullptr} on failure instead of throwing.
  \item No-throw guarantee: this function never throws exceptions.
  \item {\color{green80}\faThumbsUp}\ \ No accidental conversion.
  \item \lstinline{@TPointer@} must be a pointer type.
  \item \lstinline{@TPointer@} must be a pointer to \lstinline{const},
    if the object is \lstinline{const}.
\end{itemize}

\begin{tcolorbox}{\F{19pt}[26pt]
\begin{lstlisting}
@Tjson@ j;
std::@Tstring@* str @F=@ j.@Fget_ptr@<std::@Tstring@*>(str);
\end{lstlisting}
}\end{tcolorbox}

\pagebreak
% ------------------------------------------------------------------------------
\slidetitle{\faHandORight\,\ Best practice: Value access}

\begin{itemize}
  \item Use \lstinline{operator @TT@},
    if you are not worried about {\bfseries copying strings}
    or {\bfseries converting numbers}.
  \item Use \lstinline{@Fget_ref@<@TT@>()},
    if you are worried.
  \item Use \lstinline{@Fget_ptr@<@TT@>()},
    to avoid exceptions.
  \item Use \lstinline{@Fget_to@<@TT@>()},
    to avoid temporary values used by other \lstinline{@Fget_@}
    functions.
  \item \lstinline{@Fget@<@TT@>()} is a verbose version of
    \lstinline{operator @TT@}.
    \begin{itemize}
      \item Can be useful in \lstinline{return} statements of functions
        returning \lstinline{auto}
      \item or in assignments to types different from \lstinline{@TT@}.
    \end{itemize}
\end{itemize}

\pagebreak
% ------------------------------------------------------------------------------
\slidetitle{Conversion to and from custom types}

\begin{itemize}
  \item Specialize \lstinline{@Nnlohmann@::@Tadl_serializer@} \hfill
    \href{https://json.nlohmann.me/api/adl_serializer/}{\faFileTextO}
\end{itemize}
\vspace{-0.2em}

\begin{tcolorbox}{\F{12pt}[16pt]
\begin{lstlisting}
namespace @Nnlohmann@ {@!\vspace{0.4em}@
template <typename @TT@>
struct @Tadl_serializer@<@Nboost@::@Toptional@<@TT@>> {
    static void @Fto_json@(@Tjson@& j, const @Nboost@::@Toptional@<@TT@>& opt) {
        if (opt @F==@ @Nboost@::none) {
            j @F=@ nullptr;
        } else {
            j @F=@ *opt;
            // this will call adl_serializer<T>::to_json, which will
            // find the free function to_json in T's namespace!
        }
    }@!\vspace{0.4em}@
    static void @Ffrom_json@(const @Tjson@& j, @Nboost@::@Toptional@<@TT@>& opt) {
        if (j.@Fis_null@()) {
            opt @F=@ @Nboost@::none;
        } else {
            opt @F=@ j.@Fget@<@TT@>();
            // same as above, but with adl_serializer<T>::from_json
        }
    }
};@!\vspace{0.4em}@
}
\end{lstlisting}
}\end{tcolorbox}
\vspace{-1em}

% ------------------------------------------------------------------------------
\end{document}
